\section*{局限性}
\label{sec:limitations}

\begin{enumerate}[leftmargin=1.2em,itemsep=2pt,topsep=2pt]
\item \textbf{合成退化。}19 格网格是参数化的，而非真实采集退化。它提供剂量--响应控制，但不复现扫描仪特有的物理过程。
\item \textbf{标签来源。}NIH 标签由报告经 NLP 导出；PadChest 使用 120 标签词表并映射到八种发现，肺炎结构性缺失。标签噪声限制本文每个 AUROC 的上界。共同标签分析（附录~\ref{app:risk}）在两个数据集都标注的七个标签上重算迁移损失；排除肺炎不改变结论。
\item \textbf{跨数据集聚因。}NIH 与 PadChest 同时存在医院来源、病例构成、标签映射与输入几何的差异。共同标签分析去掉了结构性缺失的肺炎项，响应性检查去掉了常数输出解释，但这里没有任何设计能把医院偏移与病例构成偏移分开。我们主张数据集层面的判别力坍缩，而不是量化的医院域惩罚。
\item \textbf{六个模型、一块 GPU、一套解码协议。}贪婪解码最大化可复现性，但可能低估基于采样的部署所具有的输出多样性。
\item \textbf{贪婪解码不具 batch 不变性。}同一 100 张图像在固定种子下用 batch 1、4、8 重跑，有 12--13\% 的图像原始回复发生变化；解析后的状态平均在八个标签中一致 7.96--7.97 个，解析率保持 1.000。该差异比本文报告的退化效应小一个数量级，但我们不主张逐位可复现。
\item \textbf{跨家族输入几何因架构而异}（同一张图像上 0.20--15.24\,MP），无法在不改变模型的情况下拉平。我们报告它而不是隐藏它（表~\ref{tab:geometry}）；受控的输入分辨率扫描需要重跑每个模型，本文未做。
\item \textbf{没有第四个数据集。}该框架在三个数据层上展示；在独立医院或未见模型家族上的外部验证仍是未来工作。
\item \textbf{CheXagent 的契约不是同一任务。}其 79\% 的 \texttt{not\_mentioned} 率使其选择性指标不可比；我们把它排除在比较性选择性主张之外。
\item \textbf{Open-i 打分使用规则标注器，并用 CheXbert 交叉核验。}两种标注器在绝对计数与漏报排序上不一致（Spearman $-0.49$）；幻觉排序方向保持一致（$0.71$）。因此我们只把幻觉排序作为稳定结果报告，不主张普遍的漏报--幻觉权衡。
\item \textbf{监督基线的范围。}基线是单个 DenseNet-121，只用仓库内 3,871 张训练图像训练，而非完整的 11.2 万张 ChestX-ray14，因此 0.733 是这些标签所能支撑性能的下界。它足以区分``任务不可能''与``模型弱''（本文的主张），但不是临床级参照。
\item \textbf{第二个推理模型：覆盖率与预算。}Qwen3-VL-Thinking 网格在 5,120 token 预算下对 10,542 个格子中的 9,381 个给出可用答案（89.0\%），对残余做 8,192 token 抽样后为 9,401 个（89.2\%）；余下 1,141 条不终止，8,192 的 40 格抽样只能救回其中一半，而单次调用成本约翻倍。我们把不终止作为实测结果，并在共同作答的研究上报告配对召回，但逐格的非配对召回会偏向模型``愿意完成''的研究。推理链未经人工核验；我们只主张其长度与终止情况，不主张其忠实性。
\item \textbf{复现实验的覆盖范围。}schema 扫描覆盖第二个医学模型与跨家族模型（Qwen3-VL），改写对照覆盖两个模型；双轴结果仍建立在两个具有不同通道的模型上，比单模型强，但仍不是跨家族规律。CheXbert 交叉核验是标注器复现，而不是人工裁决的金标准。
\item \textbf{没有临床结局。}本文不测量诊断获益或伤害。
\item \textbf{LLaVA-Med 是实测排除。}它在 PadChest、Open-i 与一个 200 例 NIH 子集上运行，只产生两个模板且没有发现级内容（NIH 上众数占 98.7\%、结构化解析率 0\%）；它被排除在发现指标之外，但没有被排除在记录之外。Open-i 没有确定标签，因此我们无法检验它那二值的``不可判读/正常''门控是否与参考报告一致。
\item \textbf{探索性与确认性。}四次反转是带预先声明规则的描述性分析，不是预注册假设检验。所有 18 个数据层--模型单元都被报告，而不是挑选子集；迁移差值额外做了 16 单元 Bonferroni 放宽；其他处不声称 $p$ 值族错误率。
\end{enumerate}
